\section{Неоткрытые подгруппы конечного индекса}

С помощью леммы Цорна построим в $\Gal(\Qal/\Q)$ подгруппы конечного индекса, не являющиеся открытыми.

\begin{lemma}[7.28]
Пусть $V$~--- бесконечномерное векторное пространство. Для каждого $n\ge1$ существует подпространство $V_n\subset V$
с $\dim V/V_n=n$.
\end{lemma}

\begin{proof}
По лемме Цорна в $V$ есть максимальное линейно независимое подмножество; обычное рассуждение показывает, что оно порождает
$V$, т.\,е.\ является базисом. Выберем базис и возьмём в качестве $V_n$ линейную оболочку базиса без $n$ его элементов.
\end{proof}

\begin{proposition}[7.29]
Группа $\Gal(\Qal/\Q)$ имеет неоткрытые нормальные подгруппы индекса $2^n$ для всех $n\ge1$.
\end{proposition}

\begin{proof}
\emph{1) Поле $E$ и его группа.} Пусть $\ell_0=-1$, $\ell_1=2$, $\ell_2=3$, $\ell_3=5,\dots$ ($\ell_k$, $k\ge1$,~---
простые числа по порядку), и $E=\Q(\sqrt{\ell_0},\sqrt{\ell_1},\sqrt{\ell_2},\dots)\subset\CC$~--- расширение Галуа
поля $\Q$ (поле разложения многочленов $X^2-\ell_k$). Положим $E_k=\Q(\sqrt{\ell_0},\dots,\sqrt{\ell_k})$. Классы
$\ell_0,\ell_1,\dots$ линейно независимы в $\F_2$-пространстве $\Q^\times/\Q^{\times2}$ (по однозначности разложения на
простые: $\pm\ell_1^{e_1}\cdots\ell_k^{e_k}$ с $e_i\in\{0,1\}$ является квадратом рационального числа лишь при знаке $+$
и всех $e_i=0$). По теории Куммера (\rf{5.30}, \rf{5.31}(b)) $[E_k:\Q]=2^{k+1}$ и
\[
\Gal(E_k/\Q)\cong(\Z/2)^{k+1},\qquad\sigma\mapsto(s_0,\dots,s_k),\ \ \sigma(\sqrt{\ell_i})=(-1)^{s_i}\sqrt{\ell_i},
\]
причём реализуются все наборы знаков. Ограничение $E_{k+1}\to E_k$ забывает последнюю координату. По \rf{7.26a}
\[
G:=\Gal(E/\Q)\cong\varlim_k(\Z/2)^{k+1}=\prod_{i\ge0}\Z/2
\]
с топологией произведения (согласованная последовательность конечных наборов~--- это бесконечный набор).

\emph{2) Плотная подгруппа.} Пусть $H=\{(s_i)\in G\mid s_i=0\text{ для почти всех }i\}=\bigoplus_i\Z/2$. Она плотна в $G$:
любой $s=(s_i)$~--- предел последовательности <<обрезков>> $(s_0,0,0,\dots),(s_0,s_1,0,\dots),\dots$ из $H$.

\emph{3) Неоткрытые подгруппы.} Фактор $G/H$~--- векторное пространство над $\F_2$, бесконечномерное (даже несчётномерное:
$G$ несчётна, а $H$ счётна). По лемме \rf{7.28} есть подпространство индекса $2^n$, и его прообраз $G_n\subset G$~--- подгруппа
индекса $2^n$, содержащая $H$. Если бы $G_n$ была открыта, она была бы замкнута (\rf{7.1a}) и содержала бы $\bar H=G$~---
противоречие. Наконец, прообраз $G_n$ при сюръективном (\rf{7.4}) ограничении $\Gal(\Qal/\Q)\to G$~--- нормальная подгруппа
индекса $2^n$ (все подгруппы абелевой $G$ нормальны). Она не открыта: иначе она была бы замкнута и имела бы конечный
индекс, её образ $G_n$ был бы компактен, т.\,е.\ замкнут, и имел бы конечный индекс, т.\,е.\ был бы открыт.
\end{proof}

\begin{remark}[7.30]
Пусть $G=\Gal(\Qal/\Q)$. В доказательстве построена замкнутая нормальная подгруппа $N=\Gal(\Qal/E)$, для которой $G/N$~---
несчётномерное векторное пространство над $\F_2$. Каждый ненулевой линейный функционал $f\colon G/N\to\F_2$ задаёт
сюръекцию $G\to\F_2$, ядро которой~--- подгруппа индекса 2; разные $f$ дают разные подгруппы. Значит, у $G$ несчётно
много (а именно $2^{2^{\aleph_0}}$) подгрупп индекса 2. Открытых среди них лишь счётное число: они соответствуют
квадратичным расширениям $\Q$ в $\Qal$, т.\,е.\ полям $\Q(\sqrt d)$, где $d\ne1$~--- целое число, свободное от квадратов. Неоткрытая подгруппа индекса~2
плотна: её замыкание содержит её и потому имеет индекс 1 или 2, а индекс 2 означал бы, что она замкнута, а значит,
открыта (\rf{7.1a}(c)). Поэтому её неподвижное поле~--- $\Q$ (\rf{7.12}(b), \rf{7.9}).
\end{remark}

\begin{aside}[7.31]
Лемма Цорна в \rf{7.29} необходима: с ZF+DC (аксиома зависимого выбора) совместимо утверждение, что каждый гомоморфизм из
проконечной группы со счётной базой топологии в конечную группу непрерывен.
\end{aside}

\begin{aside}[7.32]
Пусть $G$~--- проконечная группа, топологически конечно порождённая (т.\,е.\ замыкание подгруппы, порождённой конечным
множеством, равно $G$). Трудная теорема, доказанная лишь в XXI веке, утверждает, что каждая
подгруппа конечного индекса в $G$ открыта; т.\,е.\ топология такой группы восстанавливается по её алгебраической
структуре. Простой частный случай~--- $\Zh$ (\rf{7.34}). Группа $\Gal(\Qal/\Q)$ не является топологически конечно порождённой:
она непрерывно отображается на $\prod_{i\ge0}\Z/2$, а образы $m$ топологических образующих порождали бы плотную подгруппу
$\prod\Z/2$; но подгруппа, порождённая $m$ элементами, имеет порядок $\le2^m$, конечна, значит, замкнута и не совпадает
с бесконечной группой $\prod\Z/2$.
\end{aside}

\begin{remark}[7.32a: итог о подгруппах конечного индекса]
В компактной группе: открытая подгруппа $\Rightarrow$ конечный индекс (\rf{7.1a}(d)); замкнутая и конечного индекса
$\Rightarrow$ открытая (\rf{7.1a}(c)); но конечный индекс сам по себе открытости не даёт (\rf{7.29}). Поэтому в бесконечной теории
Галуа конечным расширениям соответствуют \emph{открытые} подгруппы, а не подгруппы конечного индекса. Тот же пример
показывает, что из конечности орбит действия компактной группы не следует открытость стабилизаторов: $\prod\Z/2$
действует на $\{\pm1\}$ через разрывный гомоморфизм в $\Z/2$, орбита из двух точек, а стабилизатор~--- неоткрытая
подгруппа индекса 2. (Обратно, открытые стабилизаторы в компактной группе всегда дают конечные орбиты: см.\ шаг 1
доказательства \rf{7.10}.)
\end{remark}
